\chapter{Experimental Results}
\label{ch:results}

This chapter reports the accepted observations produced by the protocols in
Chapter~\ref{ch:methodology}. Results are ordered from offline model behaviour
to device execution, physical integration, communication, simulation and
external-data readiness. Each numerical result retains its experimental unit
and configuration. In particular, laptop inference calls, camera-free device
pipelines and live alert-to-decision paths are not combined into a single
latency measure.

Table~\ref{tab:results-evidence-status} states the final evidence boundary.
The distinction prevents a prepared protocol or a successful
development run from being reported as a completed experiment.

\begin{table}[H]
  \centering
  \caption{Results admitted to the final evidence set. A protocol that was not
  executed or a development-only result remains visible as a limitation rather
  than being converted into a positive finding.}
  \label{tab:results-evidence-status}
  \scriptsize
  \begin{tabularx}{\textwidth}{@{}p{0.25\textwidth}p{0.18\textwidth}X@{}}
    \toprule
    Evidence programme & Status & Result admitted in this chapter \\
    \midrule
    Six-model internal evaluation & Executed & Validation-selected operating
    points and one-time frozen-threshold test results. \\
    Conversion fidelity & Partly executed & Four YOLO ONNX comparisons, SSD
    TFLite comparison and native EfficientDet artifact readiness. \\
    Strengthened 12-cell device screen & Not executed & No model--platform ranking;
    the preparation gate remained blocked. \\
    Historical device measurements & Executed & Exact recorded pipeline and
    integrated-node measurements, each with its own configuration boundary. \\
    Concurrent-service endurance & Executed with limitations & More than 30
    minutes of live camera, inference, ROS~2, logging and flight-service
    coexistence; duration and process-survival gates passed, but Scout p95
    latency and strict Jetson swap gates failed. \\
    Stationary physical workflow & Executed, bounded & Condition-specific
    confirmation, rejection and latency observations with propellers removed. \\
    ROS~2 request--reply trials & Executed & Two distinct 300-event camera-free
    campaigns and one 1,000-event campaign under concurrent inference load. \\
    Hardware fault repetitions & Bounded subset executed & Five-repetition camera-
    interface, stale-frame, unavailable-process and late/duplicate-verdict
    outcomes; network interruption and rejection-memory tests were not executed. \\
    Mission simulation & Development evidence only & Two nominal branches, four fault
    pilots and one accepted end-to-end visual confirmed-hazard run; no formal
    success probability. \\
    External proxy evaluation & Not evaluated & Readiness observations only; no
    external predictive-performance result. \\
    \bottomrule
  \end{tabularx}
\end{table}

\section{Dataset audit and six-model predictive results}
\label{sec:results:1}

The common evaluator processed all six candidates on the frozen validation
split of 700 images and 392 annotated fire boxes. The split contained 381
background images and 307 reconstructed source-family groups. Table
\ref{tab:model-validation-results} reports the threshold selected by maximum
object-level F1 at IoU 0.50 and the metrics observed at that operating point.
The final column is the mean single-image inference-call time from the laptop
screen; it is not an edge-device or end-to-end latency.

\begin{table}[p]
  \centering
  \caption{Common validation results on 700 images and 392 fire boxes. B-FPI
  is the proportion of 381 background images containing at least one retained
  prediction. Each timing row contains 700 post-warm-up inference calls.}
  \label{tab:model-validation-results}
  \small
  \resizebox{\textwidth}{!}{%
  \begin{tabular}{@{}lrrrrrrrr@{}}
    \toprule
    Candidate & AP50--95 & AP50 & Threshold & Precision & Recall & F1 & B-FPI & Mean (ms) \\
    \midrule
    YOLOv8n-416 & 0.236 & 0.566 & 0.45 & 0.681 & 0.556 & 0.612 & 0.0026 & 9.89 \\
    YOLOv8n-640 & 0.220 & 0.542 & 0.36 & 0.553 & 0.556 & 0.555 & 0.0105 & 7.91 \\
    YOLOv8s-416 & 0.203 & 0.497 & 0.30 & 0.549 & 0.533 & 0.541 & 0.0079 & 11.14 \\
    YOLOv8s-640 & 0.213 & 0.532 & 0.22 & 0.576 & 0.548 & 0.562 & 0.0052 & 10.99 \\
    SSD MobileNet-320 & 0.185 & 0.491 & 0.81 & 0.667 & 0.500 & 0.571 & 0.0052 & 27.44 \\
    EfficientDet-320 & 0.161 & 0.461 & 0.41 & 0.582 & 0.434 & 0.497 & 0.0000 & 1263.86 \\
    \bottomrule
  \end{tabular}}
\end{table}

YOLOv8n-416 produced the largest validation AP50--95 (0.236), AP50
(0.566), precision (0.681) and F1 (0.612). YOLOv8n-416 and YOLOv8n-640 tied
for the largest recall (0.556). EfficientDet produced no above-threshold
prediction on the 381 validation background images, but it also had the lowest
validation AP50--95 and recall. The zero background count is limited to this
retained subset and threshold.

After the checkpoints, thresholds and evaluator were frozen, the one-time test
pass used 733 images, comprising 354 positive images, 379 background images
and 503 annotated fire boxes across 306 reconstructed source-family groups.
No test-driven threshold adjustment or checkpoint reselection followed.

\begin{table}[p]
  \centering
  \caption{One-time test results at validation-frozen thresholds. The grouped
  95\% F1 intervals resample the 306 retained test source-family groups.}
  \label{tab:model-test-results}
  \small
  \resizebox{\textwidth}{!}{%
  \begin{tabular}{@{}lrrrrrrrrrr@{}}
    \toprule
    Candidate & Thr. & AP50--95 & AP50 & Precision & Recall & F1 & 95\% F1 interval & TP/FP/FN & B-FPI \\
    \midrule
    YOLOv8n-416 & 0.45 & 0.119 & 0.303 & 0.541 & 0.300 & 0.386 & [0.236, 0.544] & 151/128/352 & 0.0554 \\
    YOLOv8n-640 & 0.36 & 0.101 & 0.260 & 0.409 & 0.274 & 0.329 & [0.175, 0.450] & 138/199/365 & 0.0396 \\
    YOLOv8s-416 & 0.30 & 0.130 & 0.319 & 0.508 & 0.320 & 0.393 & [0.236, 0.537] & 161/156/342 & 0.0792 \\
    YOLOv8s-640 & 0.22 & 0.133 & 0.359 & 0.545 & 0.328 & 0.409 & [0.245, 0.545] & 165/138/338 & 0.0343 \\
    SSD MobileNet-320 & 0.81 & 0.114 & 0.303 & 0.546 & 0.237 & 0.330 & [0.162, 0.479] & 119/99/384 & 0.0554 \\
    EfficientDet-320 & 0.41 & 0.168 & 0.403 & 0.679 & 0.328 & 0.442 & [0.274, 0.596] & 165/78/338 & 0.0343 \\
    \bottomrule
  \end{tabular}}
\end{table}

EfficientDet ranked first on test AP50--95, precision and F1 despite ranking
last on validation AP50--95. It tied YOLOv8s-640 for the largest test recall.
F1 decreased from validation to test for every candidate, and the grouped
intervals were broad and overlapping. Figure~\ref{fig:model-ap-validation-test}
shows the AP50--95 rank change without altering the frozen selection rule.

\begin{figure}[p]
  \centering
  \includegraphics[width=\textwidth]{model_ap_validation_test.pdf}
  \caption{AP50--95 on the 700-image validation split and the one-time
  733-image test split. Test inference used checkpoints fixed before the test
  split was opened; AP is integrated over confidence rather than evaluated at
  the selected operating threshold.}
  \label{fig:model-ap-validation-test}
\end{figure}

Only three test boxes met the COCO small-object definition at stored
resolution. A separate relative-area audit found 54 boxes occupying less than
1\% of their image. Those counts are insufficient for a general small-fire
claim and do not validate the approximately 2~cm lighter condition used later.

The two-objective validation screen placed YOLOv8n-416 and YOLOv8n-640 on the
exploratory AP--latency Pareto front (Figure~\ref{fig:validation-pareto}). The
four YOLO measurements and SSD used GPU-capable framework paths, whereas the
EfficientDet export used CPU/XNNPACK. The plot therefore describes available
runtime paths rather than an architecture-only speed comparison.

\begin{figure}[p]
  \centering
  \includegraphics[width=\textwidth]{validation_ap_latency_pareto.pdf}
  \caption[Validation AP50--95 versus laptop inference-call latency.]{Validation
  AP50--95 versus mean single-image laptop inference-call latency. Blue points
  identify the two-objective exploratory Pareto front. The logarithmic latency
  axis and runtime distinction are retained because the six candidates did not
  use an identical accelerator path.}
  \label{fig:validation-pareto}
\end{figure}

\section{Conversion fidelity and deployment outcomes}
\label{sec:results:2}

All four fixed-batch YOLO checkpoints exported to ONNX and passed the ONNX
checker. Conversion fidelity was evaluated on the same 64-image calibration
manifest used to preserve the native reference predictions. Under the frozen
match-rate and confidence-error gates, the two 416-pixel ONNX candidates
passed, while both 640-pixel ONNX candidates failed. YOLOv8n-640 missed the
0.95 match-rate gate by 0.0004; the criterion was not relaxed after inspection.

\begin{table}[H]
  \centering
  \caption{Prepared deployment artifacts and laptop conversion-fidelity
  outcomes. Ready indicates eligibility for a later runtime attempt, not
  successful execution on either edge device.}
  \label{tab:conversion-results}
  \scriptsize
  \begin{tabularx}{\textwidth}{@{}p{0.25\textwidth}p{0.16\textwidth}rrX@{}}
    \toprule
    Candidate & Prepared form & Match rate & Conf. MAE & Disposition \\
    \midrule
    YOLOv8n-416 & ONNX FP32 & 0.9603 & 0.000056 & Fidelity pass; ready for runtime attempt. \\
    YOLOv8n-640 & ONNX FP32 & 0.9496 & 0.000255 & \texttt{FIDELITY\_FAILED}. \\
    YOLOv8s-416 & ONNX FP32 & 0.9598 & 0.000223 & Fidelity pass; ready for runtime attempt. \\
    YOLOv8s-640 & ONNX FP32 & 0.9169 & 0.000212 & \texttt{FIDELITY\_FAILED}. \\
    SSD MobileNet-320 & TFLite FP32 & 0.8714 & 0.003910 & \texttt{FIDELITY\_FAILED}. \\
    EfficientDet-320 & Native TFLite & N/A & N/A & Native artifact ready for device attempt. \\
    \bottomrule
  \end{tabularx}
\end{table}

The SSD conversion allocated and produced the same number of retained outputs
as the native evaluator, but only 5,577 of 6,400 detections matched at IoU
0.50, yielding a 0.8714 match rate. Its small matched-confidence error did not
override the failed detection-match criterion. EfficientDet was already
evaluated in its native TFLite form, so cross-format fidelity was not
applicable. Its artifact was classified ready for a device attempt rather than
deployed.

The strengthened screen required 12 model--platform outcomes. In the final
evidence set,
the executor and all artifact-integrity checks had passed, but both devices
were not ready for the campaign. Consequently, none of the 12 cells contributes
latency, memory, thermal or deployment-ranking evidence in this chapter. The
historical 640-pixel Sentinel TensorRT and Scout TFLite deployments remain
separate evidence and are not invalidated by the later ONNX experiments.

\section{Exact-device runtime and endurance}
\label{sec:results:3}

The retained 3 August camera-free benchmark used the exact final deployed
artifacts, ten excluded warm-up iterations and 100 measured iterations per
device. The measured pipeline included preprocessing, runtime execution and
decoding. Table~\ref{tab:device-backend-summary} shows the descriptive results.
The platforms, models, runtimes and tensor layouts differ, so the ratio is a
system-role observation rather than a controlled model comparison.

\begin{table}[H]
  \centering
  \caption{Exact deployed backend-pipeline results, excluding camera capture,
  ROS~2 communication and decision handling (100 iterations per device).}
  \label{tab:device-backend-summary}
  \small
  \begin{tabularx}{\textwidth}{@{}p{0.21\textwidth}p{0.20\textwidth}Xrrr@{}}
    \toprule
    Platform & Runtime & Input tensor & Mean (ms) & Median (ms) & P95 (ms) \\
    \midrule
    Sentinel & TensorRT FP16 & $1\!\times\!3\!\times\!640\!\times\!640$ & 58.151 & 57.868 & 59.680 \\
    Scout & TensorFlow Lite & $1\!\times\!640\!\times\!640\!\times\!3$ & 999.431 & 991.630 & 1105.076 \\
    \bottomrule
  \end{tabularx}
\end{table}

The reciprocal mean rates were approximately 17.20 and 1.00 pipelines per
second. The Sentinel row belongs to the final engine with input
$[1,3,640,640]$, output $[1,6,8400]$ and its recorded engine hash. It is not
the historical 416-pixel engine with output $[1,6,3549]$; no value from the
latter is assigned to the final deployment.

Figure~\ref{fig:hardware-latency-ecdf} presents the complete retained
distributions together with the later communication campaign. Each panel uses
its own horizontal scale, avoiding the loss of detail that would result from
placing millisecond and approximately one-second distributions on one linear
axis.

\begin{figure}[p]
  \centering
  \includegraphics[width=\textwidth]{hardware_latency_ecdf.pdf}
  \caption{Empirical cumulative distributions for the exact deployed
  camera-free backend pipelines and the 3 August correlated ROS~2/DDS
  request--reply campaign. Dashed and dotted lines identify the median and
  95th percentile within each panel.}
  \label{fig:hardware-latency-ecdf}
\end{figure}

A separate July before-and-after software experiment measured the effect of
replacing Python anchor iteration with vectorized NumPy decoding. Under its
own fixed black-frame 5/50 protocol, mean decode time changed from 148.920 to
1.290~ms on the Sentinel path and from 83.134 to 0.992~ms on the Scout path.
Mean complete-pipeline time changed from 204.324 to 57.049~ms and from
1070.811 to 979.667~ms, respectively. These are separate serial code-revision
runs with incomplete control of background load and starting temperature.
They establish the measured software bottleneck reduction in that historical
configuration, not a confidence interval or a replacement for the 3 August
exact-deployment benchmark.

Trial \path{camera10m-threads1-20260928-r1} operated both live camera and
deployed inference paths, ROS~2 transport, asynchronous logging and 1~Hz
host-resource sampling. The periodic probe received 108 of 108 correlated
current-frame Scout verdicts; both nodes exited normally, no camera error or
logger drop was recorded, and swap did not grow. Sentinel inference covered
9,411 frames with 61.239~ms p95 latency. Scout inference covered 108 requests
with 2,504.387~ms p95 runtime. Alert-to-verdict latency was 2,529.161~ms mean
and 2,616.852~ms p95, so the run failed the registered Scout p95 target of at
most 2~s. Minimum host-wide available memory was 1,581.8~MiB on Sentinel and
4,879.5~MiB on Scout.

The observation is useful but incomplete: the Sentinel Pixhawk was unpowered,
synthetic alerts bypassed Sentinel detection, and the interval was ten minutes
rather than the registered thirty. Therefore no synchronized 30-minute
complete-stack sustainability result is inferred from this ten-minute trial
alone; the later registered run is reported below.

\begin{figure}[p]
  \centering
  \includegraphics[width=\textwidth]{camera_inference_no_pixhawk_timeseries.png}
  \caption[Ten-minute stationary camera and inference baseline.]{Ten-minute
  stationary camera and edge-inference baseline. The latency panel exposes the
  failed 2~s Scout target; the remaining panels show host-wide CPU, available
  memory and maximum reported thermal-zone temperature. The Sentinel Pixhawk
  was unpowered, so the figure is not complete-stack endurance evidence.}
  \label{fig:camera-inference-timeseries}
\end{figure}

The later registered run \path{fullstack-20260929} closed the duration and
flight-service gaps. The Sentinel resource record lasted 1,939.709~s with 1,552
samples; the Scout record lasted 1,934.836~s with 1,833 samples. Both wrappers
exited normally. The Sentinel processed 30,426 frames at 60.035~ms mean and
61.170~ms p95 inference time. Mean host CPU utilization was 32.65\% on Sentinel
and 42.44\% on Scout; minimum available memory was 1,589.9 and 4,873.8~MiB,
respectively, and the largest valid reported temperatures were 51.0 and
48.2~\(^{\circ}\)C. No logger drop was recorded. MAVProxy was active and
reported its link as operational before and after the run; a pre-existing
fence-breach message remained visible, and no vehicle command was sent.

The strict acceptance result was conditional rather than an unqualified pass.
Jetson swap increased by 0.922~MiB, violating the registered no-growth rule.
The original 360-alert probe carried the wrong trial identifier and produced
0/360 accepted requests because the Scout correctly rejected the mismatches.
That failed setup is retained. The corrected probe returned 36/36 correlated
verdicts during a separate 180.245~s interval near the end of the run. Its mean
and p95 alert-to-verdict times were 2,526.056 and 2,617.746~ms, so the registered
2~s Scout target also failed. These 36 events demonstrate live service
coexistence but are not a 30-minute continuous tiered-workflow-load result.

\begin{figure}[p]
  \centering
  \includegraphics[width=\textwidth]{full_service_resource_timeseries_20260929.png}
  \caption[Thirty-minute complete-stack resource series.]{Stationary
  complete-stack resource series for the Sentinel and Scout. Coverage exceeds
  30 minutes on both platforms. CPU, available memory, swap and maximum valid
  thermal-zone temperature are host observations at their retained sample
  times; the small Sentinel swap increase is not suppressed.}
  \label{fig:full-service-resources}
\end{figure}

\begin{figure}[p]
  \centering
  \includegraphics[width=\textwidth]{full_service_latency_distributions_20260929.png}
  \caption[Complete-stack inference and tiered-workflow latency distributions.]{The
  Sentinel distribution contains 30,426 inference frames. The Scout and
  alert-to-verdict distributions contain the 36 corrected probe events from the
  final 180.245~s and therefore describe that bounded probe, not the full
  endurance interval. The Scout p95 exceeds the registered 2~s criterion.}
  \label{fig:full-service-latency}
\end{figure}

The result partially answers RQ2: the complete stationary service configuration
survived the registered duration with stable available-memory and temperature
observations, but it did not satisfy every acceptance criterion. In particular,
the Scout latency target and strict Sentinel swap rule remain measured resource
constraints rather than discarded failures.

\section{Stationary workflow, filtering and verification}
\label{sec:results:4}

The accepted physical trials used stationary vehicles with propellers removed.
The final Sentinel ran the 640-pixel TensorRT engine and three-consecutive-frame
temporal filter. The Scout used the retained 640-pixel TFLite model with a 0.05
threshold selected after diagnostic inspection. Event identifiers matched
across alert publication, Scout receipt, Scout inference, verdict publication
and final Sentinel decision.

\begin{table}[H]
  \centering
  \caption{Accepted stationary event observations. Fractions are conditional
  counts within the recorded scene and interval; they are not population
  recall, specificity or industrial reliability estimates.}
  \label{tab:integrated-condition-results}
  \small
  \begin{tabularx}{\textwidth}{@{}p{0.28\textwidth}rrrX@{}}
    \toprule
    Condition & Requests & Confirmed & Rejected & Admissible observation \\
    \midrule
    Strong visible flame & 13 & 13 & 0 & Complete correlated path accepted every request inside the 28.740~s marked interval. \\
    Approx. 2~cm unstable flame & 10 & 1 & 9 & The same path was sensitive to the smaller, unstable condition. \\
    Forced unlit live frame & 6 & 0 & 6 & Correlated requests were rejected in the retained unchanged unlit scene. \\
    \bottomrule
  \end{tabularx}
\end{table}

The strong-flame run contained 14 total requests when one pre-interval rejected
request was included. Scout inference had mean 1256.648~ms, median
1244.841~ms and range 1214.606--1416.544~ms. Sentinel alert-to-final-decision
time had mean 1438.437~ms, median 1416.765~ms and range
1405.759--1591.737~ms. During the bounded Sentinel node run, 2,220 inference
events had mean duration 60.606~ms and maximum 174.018~ms; no logger drop was
recorded.

The small-flame run contained 11 total requests, including one pre-interval
rejection. Mean Scout inference and alert-to-final-decision times were
1278.495 and 1463.527~ms. The maximum raw Scout score was 0.088566 and only
one visible-interval request exceeded the 0.05 decision threshold. The
Sentinel again retained 2,220 inference events, with mean duration 60.537~ms,
and both loggers reported zero dropped records.

The forced unlit observations comprised one request and a five-request
replication. All six received correlated rejected verdicts. In the five-event
replication, the maximum raw fire score had mean 0.000890 and range
0.000441--0.002052; Scout inference had mean 1315.770~ms. These requests were
injected after Sentinel detection and therefore do not measure Sentinel false
alerts or environmental specificity.

The temporal filter's reset and third-frame trigger behaviour passed unit
tests. However, the accepted trials retained post-filter requests rather than
every frame-level candidate and counter transition. The number of single-frame
alerts removed, event retention and filter-added time cannot be reconstructed
without inventing unlogged observations. The planned quantitative three-stage
ablation is consequently not estimable from this campaign.

\section{Communication and recovery}
\label{sec:results:5}

Two independent camera-free request--reply campaigns each contained three
sequential 100-event trials. All 600 requests received one correlated Scout
verdict, and neither campaign recorded an event-identity mismatch or responder
duplicate. They are reported separately because their distributions and raw
records differ.

\begin{table}[H]
  \centering
  \caption{Separate bounded ROS~2/DDS application round-trip campaigns. The
  measurements use the Sentinel monotonic clock and are not one-way latency or
  packet-loss estimates.}
  \label{tab:ros-rtt-results}
  \small
  \begin{tabular}{@{}lrrrrrr@{}}
    \toprule
    Campaign & Replies & Mean & Median & P95 & P99 & Maximum \\
     &  & \multicolumn{5}{c}{milliseconds} \\
    \midrule
    29 July & 300/300 & 8.152 & 5.878 & 12.758 & 32.407 & 400.453 \\
    3 August & 300/300 & 6.046 & 5.257 & 9.457 & 24.501 & 36.120 \\
    28 September loaded & 1000/1000 & 4.897 & 4.393 & 7.984 & --- & 44.011 \\
    \bottomrule
  \end{tabular}
\end{table}

The 29 July campaign retained one 400.453~ms event below the one-second
timeout; the next three largest values were below 36~ms. This single event
raised its pooled mean and standard deviation while leaving the median near
the later campaign. The 3 August result therefore does not erase the earlier
tail. Both 300-event campaigns excluded camera acquisition, model inference,
image traffic, movement, controlled interference and long-duration load.

The 28 September campaign ran while the live D435i/TensorRT Sentinel node, the
live Scout camera/TFLite verifier, both asynchronous loggers, resource samplers
and a separate domain-85 tiered-workflow probe were active. All 1,000 events returned
one correlated response, with zero missing, mismatched, duplicate or unexpected
events. The observed delivery fraction was 1.0 (two-sided Wilson 95\% interval
0.99617--1.0). The independent load witness concurrently received 36 of 36
fresh camera-based Scout verdicts. This remains stationary application-level
evidence; the Pixhawk was unpowered, no link-layer packet capture was taken and
one-way latency is not inferred.

\begin{figure}[p]
  \centering
  \includegraphics[width=\textwidth]{loaded_ros2_rtt_distribution.png}
  \caption[Loaded 1,000-event ROS~2/DDS round-trip distribution.]{Distribution
  and empirical cumulative distribution of the 1,000 Sentinel-clock ROS~2/DDS
  round trips under concurrent live camera and inference load. The median and
  p95 markers are descriptive for this stationary SwarmNet run; one-way latency
  and link-layer loss are not inferred.}
  \label{fig:loaded-ros2-rtt}
\end{figure}

A clock-offset probe further showed why cross-host timestamps were not used as
one-way delay. The minimum-round-trip estimates placed the Sentinel and Scout
wall clocks approximately 178.535 and 201.010~ms ahead of the laptop,
respectively, with a derived Scout--Sentinel difference of approximately
22.475~ms and a conservative combined uncertainty of $\pm5.088$~ms. The
round-trip results avoid that subtraction by using only one Sentinel monotonic
clock.

Four repeated fault cases produced accepted evidence. Table
\ref{tab:hardware-fault-results} reports their specified outcomes; each case
used five unique events and retained its deliberate fault boundary.

\begin{table}[H]
  \centering
  \caption{Stationary hardware fault-injection results. Fractions describe the
  registered repetitions and are not natural failure-rate estimates.}
  \label{tab:hardware-fault-results}
  \scriptsize
  \begin{tabularx}{\textwidth}{@{}p{0.25\textwidth}p{0.25\textwidth}p{0.18\textwidth}X@{}}
    \toprule
    Fault & Expected outcome & Result & Boundary \\
    \midrule
    ROS camera interface absent & \texttt{ERROR/no\_frame} & 5/5; zero inference & Camera service itself remained active. \\
    Stream stopped after live preflight & \texttt{ERROR/stale\_frame} & 5/5; zero post-stop inference & Three normal camera/model verdicts first established readiness. \\
    Scout verifier absent & Sentinel \texttt{TIMEOUT} at 10~s & 5/5; zero Scout verdicts & Synthetic detections entered at the dispatch boundary. \\
    Two verdicts after each timeout & Reject finalized duplicates & 10/10 rejected; zero extra decisions & Deterministic late/duplicate injection, not network delay. \\
    \bottomrule
  \end{tabularx}
\end{table}

Mean fault-response latency was 47.389~ms for the unavailable camera interface
and 35.341~ms for the stale stream. Scout-process absence produced timeout
decisions at 10,023.481~ms mean. Rejection of the ten post-timeout verdicts
took 2.631~ms mean and 2.840~ms p95 at the callback. Network interruption and
repeated rejection-memory cases were not completed; consequently these results
support specific fail-closed and idempotent paths, not a general recovery-rate
or mean-time-to-recover claim.

\section{Simulation outcomes}
\label{sec:results:6}

The dual-PX4 development environment executed both nominal decision branches
through namespaced ROS~2 Jazzy and PX4 SITL interfaces. The confirmed branch
completed Sentinel patrol and hold, Scout dispatch, verification, confirmed
hold, return and landing. The rejected branch returned the Scout to standby,
resumed Sentinel patrol and ended with both vehicles landed. Table
\ref{tab:simulation-development-results} also records the retained fault
pilots.

\begin{table}[H]
  \centering
  \caption{Development simulation outcomes using provisional Gazebo airframe
  proxies and scenario-defined local-ENU targets. Each row is one retained
  pilot, not a formal repeated campaign.}
  \label{tab:simulation-development-results}
  \scriptsize
  \begin{tabularx}{\textwidth}{@{}p{0.22\textwidth}p{0.14\textwidth}p{0.16\textwidth}X@{}}
    \toprule
    Scenario & Outcome & Terminal state & Retained observation \\
    \midrule
    Confirmed hazard v1 & Pass & Both \texttt{LANDED} & Confirmed branch completed through the real SITL command interfaces. \\
    Rejected confounder v1 & Pass & Both \texttt{LANDED} & Rejected branch resumed Sentinel patrol; Scout required six mode/arm attempts after temporary health rejections. \\
    Lost alert v2 & Pass & Both \texttt{LANDED} & No Scout command attempt; typed \texttt{DEGRADED}, \texttt{RTL}, \texttt{LANDED} sequence retained. \\
    Scout unavailable v2 & Pass & Both \texttt{LANDED} & No Scout command attempt; distinct failure reason and safe sequence retained. \\
    Delayed verdict v1 & Pass & Both \texttt{LANDED} & Timeout moved both previously flown vehicles through the safe sequence. \\
    Sensor/inference error v1 & Pass & Both \texttt{LANDED} & Error retained an unverified-hazard record and initiated the safe sequence. \\
    \bottomrule
  \end{tabularx}
\end{table}

The first lost-alert and Scout-unavailable pilots exposed premature assignment
of the Scout \texttt{TAKEOFF} state despite no actuation. The controller was
corrected and the v2 logs showed a direct \texttt{PREFLIGHT} to
\texttt{DEGRADED} transition. No wrong-system command or namespace cross-talk
was observed in the retained pilots. The two nominal and four fault passes do
not estimate mission reliability. A later accepted visual confirmed-hazard run
used system ID~2 for the orange X500 V2-derived Sentinel and system ID~3 for the
blue ModalAI/Starling-derived Scout. It reached \texttt{PATROL} at 39.821~s,
received the scenario hazard at 105.215~s, began Scout transit at 133.424~s,
reached the 4.5~m commanded horizontal standoff at 144.124~s, published the
scenario-defined confirmed verdict and alarm at 147.223~s, and ended with both
vehicles landed and disarmed at 214.253~s. Figure
\ref{fig:gazebo-visual-mission} provides the visual sequence; its complete
continuous video, event log and manifest are hash-linked. The Sentinel resumed
independent patrol while the Scout verified the target. The red fire marker,
trigger, target-acquisition event and verdict were simulator scenario truth,
not outputs of a detector. The accepted v17 run used a recorded 60~s landing
bound after the retained v16 attempt missed the earlier 30~s bound during
recording; this threshold change and both attempts remain traceable.

The continuous recording from the final v17 run was subsequently inspected in
full. The qualitative review established that the retained presentation was
understandable and role-consistent for thesis and defence use;
it is not an additional simulator run and does not change the quantitative or
reliability interpretation. The v6 derived delivery preserves that video's
hash while replacing the misleading alarm still with a hash-linked event-log
panel and using the registered safe-standoff capture frame.

The preregistered 310-run campaign was not executed; several formal injectors
and the final airframe-parameter freeze remain open. The accepted visual run is
therefore one traceable development observation, not a mission-success rate.

\section{External generalization}
\label{sec:results:7}

No external predictive-performance result was admitted. FASDD-UAV licensing
was resolved to CC BY-SA 4.0 and its publisher reported 25,097 UAV samples,
but the archive had not been acquired, hashed, decoded, mapped or screened for
overlap. No FASDD-UAV file entered training or evaluation.

The MmodalFire archive and license passed integrity checks. All 28 retained
videos decoded without a decoder failure, yielding 2,974,209 frames compared
with 2,974,228 container-reported frames; the nineteen one-frame differences
remain metadata warnings. The stricter alignment consolidator validated 8 of
28 video ranges and left 20 for independent review. Event/condition alignment,
overlap checks and a final eligible manifest remained open.

These readiness observations show why RQ5 has no evaluated result rather than a
negative result.
No frozen model was run on an eligible external manifest, so the chapter
reports neither an external accuracy estimate nor a transfer conclusion.

\reviewstatus{In the final evidence set, RQ2 has a 30-minute complete-stack result with two
failed acceptance criteria; RQ3 has no common-stimulus temporal/tiered-workflow
ablation; RQ4 is supported by development simulation and one accepted visual
mission only; and RQ5 has no admitted external inference result.}
